% --- Archon physics formalization source begin ---
% archon:physics
% archon:covers PhyXMiniProblems/problem_phyx_mini_0575.lean
% archon:source-report reports/phyx_mini/problem_phyx_mini_0575.source.json

\chapter{Physics problem phyx\_mini\_0575}
\label{ch:PhyXMiniProblems_problem_phyx_mini_0575}

\paragraph{Problem source.}
\#\# Physical scenario

Doping changes the Fermi energy of a semiconductor. Consider silicon, with a gap of 1.11\textbackslash{}text\{eV\} between the top of the valence band and the bottom of the conduction band. At 300\textbackslash{}text\{K\} the Fermi level of the pure material is nearly at the mid-point of the gap. Suppose that silicon is doped with donor atoms, each of which has a state 0.15\textbackslash{}text\{eV\} below the bottom of the silicon conduction band, and suppose further that doping raises the Fermi level to 0.11\textbackslash{}text\{eV\} below the bottom of that band.

\#\# Question

Calculate the probability that a state in the doped material (at the donor level) is occupied.

\#\# Answer choices

- A: A: 0.014

- B: B: 0.176

- C: C: 0.824

- D: D: 0.986

\#\# Figure caption (auxiliary; use the image as primary evidence)

The image is a diagram illustrating energy levels related to semiconductor physics. 

- At the top, there is a blue block labeled "Conduction band."
- Below, there is a red block labeled "Valence band."
- Between the two bands, there is a vertical arrow marked with "1.11 eV," indicating the energy gap between the conduction band and valence band.
- Inside the gap, there are horizontal dashed lines representing energy levels.
- Two lines are labeled: one as "Fermi level" and the other as "Donor level."
- The Fermi level is positioned slightly above the valence band and below the donor level.
- The donor level is shown closer to the conduction band than the valence band.

\#\# Dataset metadata

Quantum Phenomena; Physical Model Grounding Reasoning, Numerical Reasoning

\paragraph{Recorded answer/context.}
C: C: 0.824

\paragraph{Figure/image path.}
/root/proposal\_for\_physic/science-mango/phyx\_mini\_run/phyx\_data/test\_image/575.png

\paragraph{Formalization target.}
create a compiling Lean file with sorry bodies at `PhyXMiniProblems/problem\_phyx\_mini\_0575.lean`. The Lean declarations must preserve the physical quantities, dimensions or dimensional roles, figure labels, governing-law hypotheses, and final relation expressed by this problem.
Use Mathlib/Physlib names found through LeanExplore where available. If a physics API is missing, introduce faithful local abstractions rather than scalar placeholder aliases.

\begin{theorem}[Physics formalization target]
\label{thm:physics:phyx_mini_0575:target}
\lean{PhyXMiniProblems.ProblemPhyXMini0575.donorLevelOccupation_matches_recordedAnswerC}
\uses{def:physics:phyx-mini-0575:phyxminiproblems-problemphyxmini0575-boltzmannconstantinelectronvoltsperkelvin, def:physics:phyx-mini-0575:phyxminiproblems-problemphyxmini0575-donordopedsemiconductorsetup, def:physics:phyx-mini-0575:phyxminiproblems-problemphyxmini0575-donorleveloccupationprobability, def:physics:phyx-mini-0575:phyxminiproblems-problemphyxmini0575-matchesdonordopedsiliconscenario, def:physics:phyx-mini-0575:phyxminiproblems-problemphyxmini0575-matchessuppliedsemiconductorenergydiagram, def:physics:phyx-mini-0575:phyxminiproblems-problemphyxmini0575-matchesproblemenergyreadouts, def:physics:phyx-mini-0575:phyxminiproblems-problemphyxmini0575-hasphysicalsemiconductorparameters, def:physics:phyx-mini-0575:phyxminiproblems-problemphyxmini0575-satisfiesfermidiracoccupationlaw, def:physics:phyx-mini-0575:phyxminiproblems-problemphyxmini0575-agreeswithinonethousandth, def:physics:phyx-mini-0575:phyxminiproblems-problemphyxmini0575-isuniquenearestdisplayedprobability}
The assigned autoformalize agent should translate this physics problem into Lean declarations in the covered file, with theorem and lemma proof bodies written as `by sorry`.
\end{theorem}
\begin{proof}
This is an autoformalization task, not a proof task. Produce statements that can later be proved by the physics prover without weakening the physical model.
\end{proof}

% --- Archon declaration topology begin ---
\section{Declaration topology}
\begin{definition}[energy In Joules]
  \label{def:physics:phyx-mini-0575:phyxminiproblems-problemphyxmini0575-energyinjoules}
  \lean{PhyXMiniProblems.ProblemPhyXMini0575.energyInJoules}
  Coherent-SI readout of an energy in joules
\end{definition}
\begin{proof}
  This is the declared model definition of energy In Joules.
\end{proof}
\begin{definition}[energy In Electron Volts]
  \label{def:physics:phyx-mini-0575:phyxminiproblems-problemphyxmini0575-energyinelectronvolts}
  \lean{PhyXMiniProblems.ProblemPhyXMini0575.energyInElectronVolts}
  \uses{def:physics:phyx-mini-0575:phyxminiproblems-problemphyxmini0575-energyinjoules}
  Read a dimensionful energy in electron-volts. The denominator is Physlib's physical one-electron-volt energy, so this is a unit readout rather than a scalar replacement for energy
\end{definition}
\begin{proof}
  This is the declared model definition of energy In Electron Volts.
\end{proof}
\begin{definition}[temperature In Kelvins]
  \label{def:physics:phyx-mini-0575:phyxminiproblems-problemphyxmini0575-temperatureinkelvins}
  \lean{PhyXMiniProblems.ProblemPhyXMini0575.temperatureInKelvins}
  Read an absolute temperature in kelvins. Physlib stores a nonnegative temperature in an arbitrary absolute scale; the unit ratio performs the calibration to TemperatureUnit.kelvin
\end{definition}
\begin{proof}
  This is the declared model definition of temperature In Kelvins.
\end{proof}
\begin{definition}[boltzmann Constant In Electron Volts Per Kelvin]
  \label{def:physics:phyx-mini-0575:phyxminiproblems-problemphyxmini0575-boltzmannconstantinelectronvoltsperkelvin}
  \lean{PhyXMiniProblems.ProblemPhyXMini0575.boltzmannConstantInElectronVoltsPerKelvin}
  \uses{def:physics:phyx-mini-0575:phyxminiproblems-problemphyxmini0575-energyinjoules}
  Boltzmann's constant expressed in electron-volts per kelvin. Physlib supplies Constants.kB in coherent SI joules per kelvin, so division by the joule readout of one electron-volt gives the required calibrated scalar
\end{definition}
\begin{proof}
  This is the declared model definition of boltzmann Constant In Electron Volts Per Kelvin.
\end{proof}
\begin{definition}[fermi Dirac Occupation Probability]
  \label{def:physics:phyx-mini-0575:phyxminiproblems-problemphyxmini0575-fermidiracoccupationprobability}
  \lean{PhyXMiniProblems.ProblemPhyXMini0575.fermiDiracOccupationProbability}
  \uses{def:physics:phyx-mini-0575:phyxminiproblems-problemphyxmini0575-boltzmannconstantinelectronvoltsperkelvin}
  The Fermi--Dirac occupation function in consistent eV and K readouts. This is a general distribution law and is not specialized to the donor level or to any displayed answer
\end{definition}
\begin{proof}
  This is the declared model definition of fermi Dirac Occupation Probability.
\end{proof}
\begin{definition}[Semiconductor Material]
  \label{def:physics:phyx-mini-0575:phyxminiproblems-problemphyxmini0575-semiconductormaterial}
  \lean{PhyXMiniProblems.ProblemPhyXMini0575.SemiconductorMaterial}
  Semiconductor material named in the problem
\end{definition}
\begin{proof}
  This is the declared model definition of Semiconductor Material.
\end{proof}
\begin{definition}[Dopant Kind]
  \label{def:physics:phyx-mini-0575:phyxminiproblems-problemphyxmini0575-dopantkind}
  \lean{PhyXMiniProblems.ProblemPhyXMini0575.DopantKind}
  Kind of impurity introduced into the semiconductor
\end{definition}
\begin{proof}
  This is the declared model definition of Dopant Kind.
\end{proof}
\begin{definition}[Energy Reference]
  \label{def:physics:phyx-mini-0575:phyxminiproblems-problemphyxmini0575-energyreference}
  \lean{PhyXMiniProblems.ProblemPhyXMini0575.EnergyReference}
  Physical energy references appearing in the prose and diagram
\end{definition}
\begin{proof}
  This is the declared model definition of Energy Reference.
\end{proof}
\begin{definition}[Diagram Feature]
  \label{def:physics:phyx-mini-0575:phyxminiproblems-problemphyxmini0575-diagramfeature}
  \lean{PhyXMiniProblems.ProblemPhyXMini0575.DiagramFeature}
  Features explicitly visible in the supplied energy-level diagram
\end{definition}
\begin{proof}
  This is the declared model definition of Diagram Feature.
\end{proof}
\begin{definition}[Diagram Mark Style]
  \label{def:physics:phyx-mini-0575:phyxminiproblems-problemphyxmini0575-diagrammarkstyle}
  \lean{PhyXMiniProblems.ProblemPhyXMini0575.DiagramMarkStyle}
  Rendering style used for a feature in the supplied diagram
\end{definition}
\begin{proof}
  This is the declared model definition of Diagram Mark Style.
\end{proof}
\begin{definition}[Semiconductor Energy Diagram]
  \label{def:physics:phyx-mini-0575:phyxminiproblems-problemphyxmini0575-semiconductorenergydiagram}
  \lean{PhyXMiniProblems.ProblemPhyXMini0575.SemiconductorEnergyDiagram}
  \uses{def:physics:phyx-mini-0575:phyxminiproblems-problemphyxmini0575-diagramfeature, def:physics:phyx-mini-0575:phyxminiproblems-problemphyxmini0575-diagrammarkstyle}
  Primary-image data. Vertical coordinates are schematic dimensionless drawing coordinates increasing upward; they are not physical energies. The numerical gap annotation is an electron-volt readout printed beside the arrow
\end{definition}
\begin{proof}
  This is the declared model definition of Semiconductor Energy Diagram.
\end{proof}
\begin{definition}[Donor Doped Semiconductor Setup]
  \label{def:physics:phyx-mini-0575:phyxminiproblems-problemphyxmini0575-donordopedsemiconductorsetup}
  \lean{PhyXMiniProblems.ProblemPhyXMini0575.DonorDopedSemiconductorSetup}
  \uses{def:physics:phyx-mini-0575:phyxminiproblems-problemphyxmini0575-semiconductormaterial, def:physics:phyx-mini-0575:phyxminiproblems-problemphyxmini0575-dopantkind, def:physics:phyx-mini-0575:phyxminiproblems-problemphyxmini0575-energyreference, def:physics:phyx-mini-0575:phyxminiproblems-problemphyxmini0575-semiconductorenergydiagram}
  Independent physical data for the pure and donor-doped silicon samples. The energy function retains all levels as dimensionful quantities. The tolerance parameter gives quantitative meaning to the prose word "nearly" without inventing a fixed numerical tolerance not stated by the source
\end{definition}
\begin{proof}
  This is the declared model definition of Donor Doped Semiconductor Setup.
\end{proof}
\begin{definition}[donor Level Occupation Probability]
  \label{def:physics:phyx-mini-0575:phyxminiproblems-problemphyxmini0575-donorleveloccupationprobability}
  \lean{PhyXMiniProblems.ProblemPhyXMini0575.donorLevelOccupationProbability}
  \uses{def:physics:phyx-mini-0575:phyxminiproblems-problemphyxmini0575-donordopedsemiconductorsetup}
  The dimensionless probability requested by the question
\end{definition}
\begin{proof}
  This is the declared model definition of donor Level Occupation Probability.
\end{proof}
\begin{definition}[Is Near Band Gap Midpoint Within]
  \label{def:physics:phyx-mini-0575:phyxminiproblems-problemphyxmini0575-isnearbandgapmidpointwithin}
  \lean{PhyXMiniProblems.ProblemPhyXMini0575.IsNearBandGapMidpointWithin}
  \uses{def:physics:phyx-mini-0575:phyxminiproblems-problemphyxmini0575-energyinelectronvolts}
  An energy is near the midpoint of two band edges when its electron-volt readout lies within a nonnegative tolerance smaller than half the gap. The tolerance remains scenario data because the source says only "nearly"
\end{definition}
\begin{proof}
  This is the declared model definition of Is Near Band Gap Midpoint Within.
\end{proof}
\begin{definition}[Matches Donor Doped Silicon Scenario]
  \label{def:physics:phyx-mini-0575:phyxminiproblems-problemphyxmini0575-matchesdonordopedsiliconscenario}
  \lean{PhyXMiniProblems.ProblemPhyXMini0575.MatchesDonorDopedSiliconScenario}
  \uses{def:physics:phyx-mini-0575:phyxminiproblems-problemphyxmini0575-energyinelectronvolts, def:physics:phyx-mini-0575:phyxminiproblems-problemphyxmini0575-temperatureinkelvins, def:physics:phyx-mini-0575:phyxminiproblems-problemphyxmini0575-donordopedsemiconductorsetup, def:physics:phyx-mini-0575:phyxminiproblems-problemphyxmini0575-isnearbandgapmidpointwithin}
  The material, donor doping, temperature, qualitative intrinsic-level placement, and the statement that doping raises the Fermi level. None of these fields asserts the requested occupation probability
\end{definition}
\begin{proof}
  This is the declared model definition of Matches Donor Doped Silicon Scenario.
\end{proof}
\begin{definition}[Matches Supplied Semiconductor Energy Diagram]
  \label{def:physics:phyx-mini-0575:phyxminiproblems-problemphyxmini0575-matchessuppliedsemiconductorenergydiagram}
  \lean{PhyXMiniProblems.ProblemPhyXMini0575.MatchesSuppliedSemiconductorEnergyDiagram}
  \uses{def:physics:phyx-mini-0575:phyxminiproblems-problemphyxmini0575-semiconductorenergydiagram}
  Labels, drawing styles, gap annotation, and bottom-to-top order read from the primary image 575.png. The bitmap places the donor dashed line below the Fermi dashed line, consistently with the stated 0.15 eV and 0.11 eV depths; this takes precedence over the contradictory auxiliary caption
\end{definition}
\begin{proof}
  This is the declared model definition of Matches Supplied Semiconductor Energy Diagram.
\end{proof}
\begin{definition}[Matches Problem Energy Readouts]
  \label{def:physics:phyx-mini-0575:phyxminiproblems-problemphyxmini0575-matchesproblemenergyreadouts}
  \lean{PhyXMiniProblems.ProblemPhyXMini0575.MatchesProblemEnergyReadouts}
  \uses{def:physics:phyx-mini-0575:phyxminiproblems-problemphyxmini0575-energyinelectronvolts, def:physics:phyx-mini-0575:phyxminiproblems-problemphyxmini0575-donordopedsemiconductorsetup}
  The calibrated physical energy differences stated in the problem. They relate independent dimensionful levels and do not mention an occupation probability or answer choice
\end{definition}
\begin{proof}
  This is the declared model definition of Matches Problem Energy Readouts.
\end{proof}
\begin{definition}[Has Physical Semiconductor Parameters]
  \label{def:physics:phyx-mini-0575:phyxminiproblems-problemphyxmini0575-hasphysicalsemiconductorparameters}
  \lean{PhyXMiniProblems.ProblemPhyXMini0575.HasPhysicalSemiconductorParameters}
  \uses{def:physics:phyx-mini-0575:phyxminiproblems-problemphyxmini0575-energyinelectronvolts, def:physics:phyx-mini-0575:phyxminiproblems-problemphyxmini0575-temperatureinkelvins, def:physics:phyx-mini-0575:phyxminiproblems-problemphyxmini0575-donordopedsemiconductorsetup}
  Positivity, in-gap placement, and probability bounds selecting the physical branch of the model. These conditions contain no numerical value for the donor occupation probability
\end{definition}
\begin{proof}
  This is the declared model definition of Has Physical Semiconductor Parameters.
\end{proof}
\begin{definition}[Satisfies Fermi Dirac Occupation Law]
  \label{def:physics:phyx-mini-0575:phyxminiproblems-problemphyxmini0575-satisfiesfermidiracoccupationlaw}
  \lean{PhyXMiniProblems.ProblemPhyXMini0575.SatisfiesFermiDiracOccupationLaw}
  \uses{def:physics:phyx-mini-0575:phyxminiproblems-problemphyxmini0575-energyinelectronvolts, def:physics:phyx-mini-0575:phyxminiproblems-problemphyxmini0575-temperatureinkelvins, def:physics:phyx-mini-0575:phyxminiproblems-problemphyxmini0575-fermidiracoccupationprobability, def:physics:phyx-mini-0575:phyxminiproblems-problemphyxmini0575-donordopedsemiconductorsetup}
  The governing Fermi--Dirac law for every electronic-state energy in the doped sample. It relates the independent occupation observable to the doped Fermi level and temperature generally; it is not specialized to the donor state or to the recorded answer
\end{definition}
\begin{proof}
  This is the declared model definition of Satisfies Fermi Dirac Occupation Law.
\end{proof}
\begin{lemma}[donor Energy Relative To Doped Fermi]
  \label{lem:physics:phyx-mini-0575:phyxminiproblems-problemphyxmini0575-donorenergyrelativetodopedfermi}
  \lean{PhyXMiniProblems.ProblemPhyXMini0575.donorEnergyRelativeToDopedFermi}
  \uses{def:physics:phyx-mini-0575:phyxminiproblems-problemphyxmini0575-energyinelectronvolts, def:physics:phyx-mini-0575:phyxminiproblems-problemphyxmini0575-donordopedsemiconductorsetup, def:physics:phyx-mini-0575:phyxminiproblems-problemphyxmini0575-matchesproblemenergyreadouts}
  Subtracting the two stated depths below the conduction-band edge places the donor state 0.04 eV below the doped Fermi level. This is a derived intermediate relation, not a premise
\end{lemma}
\begin{proof}
  The conclusion follows from the governing relations and figure readouts named in the statement of donor Energy Relative To Doped Fermi.
\end{proof}
\begin{definition}[Answer Choice]
  \label{def:physics:phyx-mini-0575:phyxminiproblems-problemphyxmini0575-answerchoice}
  \lean{PhyXMiniProblems.ProblemPhyXMini0575.AnswerChoice}
  Labels attached to the four displayed occupation probabilities
\end{definition}
\begin{proof}
  This is the declared model definition of Answer Choice.
\end{proof}
\begin{definition}[displayed Occupation Probability]
  \label{def:physics:phyx-mini-0575:phyxminiproblems-problemphyxmini0575-displayedoccupationprobability}
  \lean{PhyXMiniProblems.ProblemPhyXMini0575.displayedOccupationProbability}
  \uses{def:physics:phyx-mini-0575:phyxminiproblems-problemphyxmini0575-answerchoice}
  Dimensionless probability printed beside each answer label
\end{definition}
\begin{proof}
  This is the declared model definition of displayed Occupation Probability.
\end{proof}
\begin{definition}[recorded Dataset Answer]
  \label{def:physics:phyx-mini-0575:phyxminiproblems-problemphyxmini0575-recordeddatasetanswer}
  \lean{PhyXMiniProblems.ProblemPhyXMini0575.recordedDatasetAnswer}
  \uses{def:physics:phyx-mini-0575:phyxminiproblems-problemphyxmini0575-answerchoice}
  Dataset metadata recording answer C; this is not a theorem premise
\end{definition}
\begin{proof}
  This is the declared model definition of recorded Dataset Answer.
\end{proof}
\begin{definition}[Agrees Within One Thousandth]
  \label{def:physics:phyx-mini-0575:phyxminiproblems-problemphyxmini0575-agreeswithinonethousandth}
  \lean{PhyXMiniProblems.ProblemPhyXMini0575.AgreesWithinOneThousandth}
  \uses{def:physics:phyx-mini-0575:phyxminiproblems-problemphyxmini0575-answerchoice, def:physics:phyx-mini-0575:phyxminiproblems-problemphyxmini0575-displayedoccupationprobability}
  The exact standard-constant calculation differs from the displayed 0.824 by less than 0.001. This tolerance describes agreement with the supplied three-decimal answer without falsely claiming exact equality or conventional rounding at the half-thousandth boundary
\end{definition}
\begin{proof}
  This is the declared model definition of Agrees Within One Thousandth.
\end{proof}
\begin{definition}[Is Unique Nearest Displayed Probability]
  \label{def:physics:phyx-mini-0575:phyxminiproblems-problemphyxmini0575-isuniquenearestdisplayedprobability}
  \lean{PhyXMiniProblems.ProblemPhyXMini0575.IsUniqueNearestDisplayedProbability}
  \uses{def:physics:phyx-mini-0575:phyxminiproblems-problemphyxmini0575-answerchoice, def:physics:phyx-mini-0575:phyxminiproblems-problemphyxmini0575-displayedoccupationprobability}
  A choice is strictly closer than every distinct displayed alternative
\end{definition}
\begin{proof}
  This is the declared model definition of Is Unique Nearest Displayed Probability.
\end{proof}
% --- Archon declaration topology end ---
% --- Archon physics formalization source end ---
